% release notes for GRIDGEN version 1.0.03
\item Version 1.0.03 10/31/2026
\begin{itemize}[label=--]
  \item[] \textbf{BUG FIXES AND OTHER CHANGES TO EXISTING FUNCTIONALITY}
  \item Layer records in QUADTREE\_BUILDER and QUADTREE blocks were ignored, without a message, after the first layer that lacked them: an ACTIVE\_DOMAIN LAYER record unless every layer above it had one, and refinement feature, refinement array, TOP LAYER, and BOTTOM LAYER records after the first layer without a BOTTOM LAYER record. Each layer's records are now read independently, so a grid that names records for a lower layer only can change.
  \item The record names REFINE\_LEVEL\_BY\_ATTRIBUTE, HORIZONTAL\_SMOOTHING\_LEVEL, and VERTICAL\_SMOOTHING\_LEVEL given in the GRIDGEN report are now accepted. Previously, the first stopped GRIDGEN and the others were ignored, so smoothing stayed at level 1. The spellings GRIDGEN accepted before, REFINEMENT\_LEVEL\_BY\_ATTRIBUTE, SMOOTHING\_LEVEL\_HORIZONTAL, and SMOOTHING\_LEVEL\_VERTICAL, are still accepted, and the report's spelling is used if both are given. A refinement level read from a shapefile attribute now overrides REFINEMENT\_LEVEL, as the report describes; previously REFINEMENT\_LEVEL was used when both were given. GRIDGEN now prints a warning for each record it does not recognize in a block, such as a misspelled name, instead of ignoring it.
  \item Cell elevations read with the ASCIIGRID option and the default area weighting were averaged over a four-sided area whose upper edge extended half the cell height, rather than half the cell width, east of the cell center. For a cell that is not square, the average therefore included part of the ASCII grid beyond the cell's upper-right corner, or left out part of the cell (for example, 5.7 instead of 6.0 for a cell 10 wide and 20 tall). Each cell is now averaged over its own area. Square cells, and ASCII grids that do not vary across a cell, are unchanged.
  \item Cell elevations read with the ASCIIGRID option included NODATA values in the area-weighted average, so a cell partly over missing data was given an elevation pulled toward the NODATA value (for example, about -2,500 instead of -42 for a cell one quarter over a -9999 value). Missing data are now left out of the average. With either the area-weighted or the bilinear interpolation, a cell with no valid value to interpolate from (for example, one that lies entirely over NODATA values) now stops GRIDGEN with an error naming the ASCII grid and the cell location, instead of being given the NODATA value as its elevation. Elevations from ASCII grids without NODATA values are unchanged, and a cell beyond the extent of an ASCII grid still takes the value at its edge.
  \item A polygon that enclosed the whole base grid, without its boundary crossing the grid, was treated as if it covered no cells. As an ACTIVE\_DOMAIN it made every cell inactive; every cell is now active, and the grid is the same as one built without an active domain. As a REFINEMENT\_FEATURES polygon it refined nothing; every cell in the layer is now refined to the REFINEMENT\_LEVEL, as the cells inside a polygon that crosses the grid already were. In a GRID\_INTERSECTION block it selected nothing; it now selects every cell in the layer. A polygon that lies outside the grid still covers no cells.
  \item ROTATION\_ANGLE was converted between degrees and radians with pi taken as 3.1415926, so a rotated grid differed from the same grid rotated with an exact pi (for example, by FloPy) by about 0.6 millimeters at 110 kilometers from the grid origin. The conversion now uses pi to double precision, which moves the cells of a rotated grid and changes intersection results slightly. The grid definition file written by QUADTREE\_BUILDER now repeats the rotation angle as given (17.3 rather than 17.30000000000000071054273576010018587112).
  \item A quadtree grid built with NODE\_NUMBERING = zero-based wrote a tree structure file that later actions could not read. The file now gives each cell's layer, row, and column starting at 1 and its quadrant numbers from 1 to 4, as the GRIDGEN report describes, so it can be read back, and a grid read from it keeps zero-based node numbers in all exported files. A zero-based tree structure file written by an earlier version still stops GRIDGEN; rebuild the grid. Output for grids built with the default one-based numbering is unchanged.
  \item Real values in several output files were written with too few digits. VTK files with shared vertices stored coordinates and cell values in single precision, and all VTK files wrote six significant digits (2.86978e+06 for a northing, or 10-meter resolution); connection lengths and areas (the c1, c2, and fahl files) were stored in single precision; array files were written with 10 significant digits; and ghost-node coefficients and intersection (.ifo) values were written with six. The .nod file and the grid definition file used 40 digits, which writes digits that do not represent the stored value. Each real value is now stored in double precision and written with the fewest digits, up to 17, that read back to the same value.
  \item Each cell's corners were computed from its own center and size, so a corner shared by adjacent cells could be written with coordinates that differ in their last digits when cell sizes are not exactly representable in binary and coordinates are large (for example, UTM or state plane). Shared vertices in grid shapefiles and VTK files therefore had to be matched with a tolerance, and rounding did not merge all of them. Each corner is now computed once, and every cell that shares it is written with the same coordinates. For a rotated test grid of 812.7 by 793.3 cells, the polygon shapefile now has 107,409 distinct corners, one for each corner of the grid, rather than 163,390.
  \item GRID\_TO\_USGDATA now writes correct output for a MODFLOW\_GRID as well as a QUADTREE grid; the GRIDGEN report describes it only for a QUADTREE grid. For a MODFLOW\_GRID, the horizontal connection areas (fahl) were wrong wherever a layer was not a whole number of length units thick or neighboring cells differed in thickness. Upward connections were also marked as downward in the connection directions (fldr), the last value of the index array (ia) was one too small, and GRIDGEN stopped before writing the node (.nod) file. VERTICAL\_PASS\_THROUGH applies only to a QUADTREE grid; for a MODFLOW\_GRID it is ignored with a warning. Output for quadtree grids is unchanged.
  \item Vertex elevations in shared-vertex VTK files (SHARE\_VERTEX = TRUE) could be incorrect where three active cells of different sizes meet, such as next to an inactive cell, and could change from run to run. They are now interpolated from the correct cells, and the files are the same on every run.
  \item Polygon intersection output listed cells in an order that could change from run to run. Cells are now listed in the same order on every run, beginning with the cells along the polygon boundary from its first vertex. The cells and values are unchanged.
  \item[] \textbf{BUILD, TESTING, AND INTERNAL CHANGES}
  \item GRIDGEN now has GoogleTest unit tests and regression tests that run 22 cases on every compiler in the continuous integration matrix. Every run checks that two runs of the same build give identical output, and pull requests also compare every output file byte for byte with the base branch. Defects found by the new tests are tracked as GitHub issues.
  \item Executables for macOS on Intel processors are no longer provided; on those computers, build GRIDGEN from the source code. The macOS executable (macarm.zip) runs on Apple silicon with macOS 14 or later.
  \item The version is set in \texttt{version.txt} and shown by the program, which adds the git commit to the version of a development build (for example, 1.0.03.dev0+d418a91). \texttt{pixi run update-version} updates the version in every file that holds it. The Visual Studio project files were removed; GRIDGEN is built with meson.
  \item[] \textbf{DOCUMENTATION}
  \item The release notes describe how GRIDGEN output defines MODFLOW 6 DISU and DISV grids and how FloPy uses GRIDGEN output to create MODFLOW-USG and MODFLOW 6 input.
\end{itemize}
